\documentclass[10pt]{article}

\usepackage[letterpaper,portrait,margin=0.22in]{geometry}
\usepackage{amsmath,amssymb,mathtools}
\usepackage{microtype}
\usepackage{multicol}
\usepackage{enumitem}
\usepackage[T1]{fontenc}
\usepackage{lmodern}

\setlength{\parindent}{0pt}
\setlength{\parskip}{1pt}
\setlength{\columnsep}{0.12in}
\setlength{\columnseprule}{0pt}
\setlength{\abovedisplayskip}{2.5pt}
\setlength{\belowdisplayskip}{2.5pt}
\setlength{\abovedisplayshortskip}{1.5pt}
\setlength{\belowdisplayshortskip}{1.5pt}
\setlist[enumerate]{leftmargin=*,nosep,label=(\alph*)}
\pagestyle{plain}
\renewcommand{\arraystretch}{1.05}

\newcommand{\examlabel}[1]{%
  \par\noindent\textbf{#1}\par\vspace{1pt}}
\newcommand{\question}[1]{\par\noindent\textbf{Question.} #1\par}
\newcommand{\answer}{\par\noindent\textbf{Answer.} }
\newcommand{\conclusion}[1]{%
  \par\noindent\framebox[\columnwidth][l]{%
    \parbox{\dimexpr\columnwidth-2\fboxsep-2\fboxrule}{\textbf{Conclusion.} #1}}\par}

\begin{document}
\fontsize{8.2}{9.45}\selectfont

% PAGE 1 ----------------------------------------------------------------------
\begin{multicols*}{3}

\examlabel{Exam 1 --- Question 2}
\question{Apply gradient descent to
\[
  \begin{bmatrix}1&3\\3&1\end{bmatrix}
  \begin{bmatrix}x\\y\end{bmatrix}
  =
  \begin{bmatrix}1\\1\end{bmatrix}.
\]
The matrix has eigenvalues $\lambda_1=4$ and $\lambda_2=-2$, with
eigenvectors
$v_1=\frac1{\sqrt2}(1,1)$ and $v_2=\frac1{\sqrt2}(1,-1)$.
The unique solution is $x^\star=(1/4,1/4)$.}

\textbf{(a)} What is the maximum rate $\rho^\star$ below which convergence
is guaranteed for every initialization?

\answer Let $A$ denote the matrix and use the sum-of-squares loss
\[
  F(z)=\lVert Az-b\rVert_2^2,
  \qquad
  \nabla F(z)=2A^\mathsf{T}(Az-b).
\]
If $e_t=z_t-z^\star$, then one gradient-descent step gives
\[
  e_{t+1}=(I-2\rho A^\mathsf{T}A)e_t.
\]
Since $A$ is symmetric, an eigenvalue $\lambda$ of $A$ produces eigenvalue
$\lambda^2$ for $A^\mathsf{T}A$.  The corresponding error multiplier is
$1-2\rho\lambda^2$.  Convergence for every initialization requires
\[
  |1-2\rho\lambda^2|<1
  \quad\text{for both eigenvalues.}
\]
The largest squared eigenvalue is $4^2=16$, so
\[
  0<\rho<\frac1{16},
  \qquad
  \boxed{\rho^\star=\frac1{16}}.
\]

\textbf{(b)} With $\rho=0.01$ and $z_0=(1,0)$, the distance to $z^\star$
is $\Theta(b^t)$. Find $b$.

\answer The two error multipliers are
\[
  q_1=1-2(0.01)(4^2)=0.68,
  \qquad
  q_2=1-2(0.01)(-2)^2=0.92.
\]
The initial error has a nonzero component in the $v_2$ direction.
Therefore the slower surviving factor controls the asymptotic distance:
\[
  \lVert z_t-z^\star\rVert_2=\Theta(0.92^t),
  \qquad
  \boxed{b=0.92}.
\]

\examlabel{Exam 1 --- Question 3}
\question{Calculate the Boolean Fourier transforms below.  The input order
is $(++,-+,+-,--)$.}

\textbf{(a)} $f=[1,1,-1,-1]$.

\answer This is exactly the function $f(x_1,x_2)=x_2$.  Hence
\[
  \widehat f(\{2\})=1,
  \qquad
  \widehat f(S)=0\quad(S\neq\{2\}).
\]

\textbf{(b)} $g=[0,1,1,1]$.

\answer The function $1-g$ is the indicator of the point $(+,+)$:
\[
  1-g
  =\operatorname{point}_{++}(x_1,x_2)
  =\frac{(1+x_1)(1+x_2)}4.
\]
Expanding and subtracting from $1$ gives
\[
  \boxed{
  g=\frac34-\frac14x_1-\frac14x_2-\frac14x_1x_2}.
\]
Thus its coefficients at
$S=\varnothing,\{1\},\{2\},\{1,2\}$ are
$3/4,-1/4,-1/4,-1/4$.

\textbf{(c)} Let
\[
  h(x_1,x_2,x_3)=
  \begin{cases}
    f(x_1,x_2),&x_3=-1,\\
    g(x_1,x_2),&x_3=+1.
  \end{cases}
\]

\answer The factors $(1-x_3)/2$ and $(1+x_3)/2$ select the two cases:
\[
  h=\frac{1-x_3}{2}f+\frac{1+x_3}{2}g
   =\frac{f+g}{2}+x_3\frac{g-f}{2}.
\]
Substituting $f=x_2$ and the polynomial for $g$ gives
\[
\boxed{
\begin{aligned}
h={}&\frac38-\frac18x_1+\frac38x_2-\frac18x_1x_2\\
 &+\frac38x_3-\frac18x_1x_3-\frac58x_2x_3
   -\frac18x_1x_2x_3 .
\end{aligned}}
\]
Therefore, in the order
$\varnothing,1,2,3,12,13,23,123$, the coefficient list is
\[
  \boxed{\frac38,-\frac18,\frac38,\frac38,
  -\frac18,-\frac18,-\frac58,-\frac18}.
\]

\examlabel{Exam 2 --- Question 2}
\question{Let $A$ be the adjacency matrix of the 3-cycle, with zeros on the
diagonal. Its largest eigenvalue is $\lambda_1=2$, with eigenvector
$v_1=\frac1{\sqrt3}(1,1,1)$.}

\textbf{(a)} If $x\perp v_1$, prove that normalized power iteration
oscillates between $x$ and $-x$.

\answer Write $x=(x_1,x_2,x_3)$. Orthogonality gives
\[
  x\cdot v_1=0
  \quad\Longrightarrow\quad
  x_1+x_2+x_3=0.
\]
Multiplication by the adjacency matrix yields
\[
  Ax=(x_2+x_3,\;x_3+x_1,\;x_1+x_2).
\]
Because the coordinate sum is zero,
$x_2+x_3=-x_1$, $x_3+x_1=-x_2$, and
$x_1+x_2=-x_3$. Therefore
\[
  Ax=(-x_1,-x_2,-x_3)=-x.
\]
If $x$ is normalized, then $\lVert Ax\rVert=\lVert-x\rVert=1$.
Thus one iteration maps $x$ to $-x$ and the next maps $-x$ back to $x$.
\[
  \boxed{x,-x,x,-x,\ldots}
\]

\textbf{(b)} Deduce the other two eigenvalues.

\answer The orthogonal complement of $v_1$ is two-dimensional, and the
calculation above shows that $A$ acts as $-I$ on every vector in that
space. Hence
\[
  \boxed{\lambda_2=\lambda_3=-1}.
\]

\examlabel{Exam 2 --- Question 3}
\question{At $x=1,i,-1,-i$, the values of $f$ are $0,1,2,3$.
Use
\[
  f(x)=f^+(x^2)+x f^-(x^2)
\]
to recover the polynomial representation of $f$.}

\textbf{(a)} Find the value lists of $f^+$ and $f^-$.

\answer Evaluate the decomposition at $x$ and $-x$. Adding eliminates the
odd part and subtracting eliminates the even part:
\[
  f^+(x^2)=\frac{f(x)+f(-x)}2,
  \qquad
  f^-(x^2)=\frac{f(x)-f(-x)}{2x}.
\]
At $x=1$,
\[
  f^+(1)=\frac{0+2}{2}=1,
  \qquad
  f^-(1)=\frac{0-2}{2}=-1.
\]
At $x=i$, where $i^2=-1$,
\[
  f^+(-1)=\frac{1+3}{2}=2,
  \qquad
  f^-(-1)=\frac{1-3}{2i}=-\frac1i=i.
\]
Thus, over the domain $(1,-1)$,
\[
  \boxed{f^+=[1,2]},
  \qquad
  \boxed{f^-=[-1,i]}.
\]

\textbf{(b)} Find the degree-one polynomial representations of $f^+$ and
$f^-$.

\answer If $p(y)=\alpha+\beta y$, $p(1)=u$, and $p(-1)=v$, then
\[
  \alpha=\frac{u+v}{2},
  \qquad
  \beta=\frac{u-v}{2}.
\]
Applying this formula to the two lists gives
\[
  f^+(y)=\frac32-\frac12y,
\]
\[
  f^-(y)=\frac{-1+i}{2}+\frac{-1-i}{2}y.
\]

\textbf{(c)} Reconstruct $f$.

\answer Substitute $y=x^2$ and multiply the odd polynomial by $x$:
\[
\boxed{
  f(x)=\frac32+\frac{-1+i}{2}x-\frac12x^2
       +\frac{-1-i}{2}x^3}.
\]

\end{multicols*}
\newpage

% PAGE 2 ----------------------------------------------------------------------
\begin{multicols*}{3}

\examlabel{Exam 3 --- Question 1}
\question{Apply gradient descent to the inconsistent scalar system
$x=0$ and $x=1$.}

\textbf{(a)} Write the sum-of-squares loss and find its minimizer.

\answer The residuals are $x$ and $x-1$, so
\[
  F(x)=x^2+(x-1)^2.
\]
Differentiate:
\[
  F'(x)=2x+2(x-1)=4x-2.
\]
The minimum occurs where $F'(x^\star)=0$, hence
\[
  \boxed{x^\star=\frac12}.
\]

\textbf{(b)} Find
$\dfrac{x_t-x^\star}{x_{t-1}-x^\star}$.

\answer One gradient-descent step is
\[
  x_t=x_{t-1}-\rho(4x_{t-1}-2).
\]
Subtract $x^\star=1/2$ and factor:
\[
  x_t-\frac12
  =(1-4\rho)\left(x_{t-1}-\frac12\right).
\]
Therefore
\[
  \boxed{\frac{x_t-x^\star}{x_{t-1}-x^\star}=1-4\rho}.
\]

\textbf{(c)} If $x_0=0$ and $\rho=1/8$, find $x_{20}$.

\answer The multiplier is $1-4(1/8)=1/2$. Applying the recurrence twenty
times gives
\[
  x_{20}-\frac12
  =\left(\frac12\right)^{20}
   \left(0-\frac12\right)
  =-\frac1{2^{21}}.
\]
Thus
\[
  \boxed{x_{20}=\frac12-\frac1{2^{21}}}.
\]

\examlabel{Exam 3 --- Question 2(a--b)}
\question{Apply subspace iteration to
\[
 A=
 \begin{bmatrix}
 0&1&2\\
 1&0&0\\
 2&0&0
 \end{bmatrix},
\]
starting from the standard basis $e_1,e_2,e_3$. Ignore $e_3$ initially.}

\textbf{(a)} Show the first three rounds for the first two vectors.

\answer Define
\[
  u=\frac1{\sqrt5}(0,1,2).
\]
\emph{Round 1.}
\[
  Ae_1=(0,1,2)=\sqrt5\,u,
  \qquad
  Ae_2=(1,0,0)=e_1.
\]
After normalization, $(x_1,x_2)=(u,e_1)$.

\emph{Round 2.}
\[
  Au=(\sqrt5,0,0)=\sqrt5\,e_1,
  \qquad
  Ae_1=\sqrt5\,u.
\]
After normalization, $(x_1,x_2)=(e_1,u)$.

\emph{Round 3.} The same computation switches them again, so
$(x_1,x_2)=(u,e_1)$.

\textbf{(b)} Use this behavior to find two eigenvalues.

\answer Since $Ae_1=\sqrt5\,u$ and $Au=\sqrt5\,e_1$,
\[
  A(e_1+u)=\sqrt5(e_1+u),
\]
\[
  A(e_1-u)=-\sqrt5(e_1-u).
\]
Therefore
\[
  \boxed{\lambda_1=\sqrt5,\qquad \lambda_2=-\sqrt5}.
\]

\examlabel{Exam 3 --- Question 2(c--d)}
\question{Prove that if $Ax_i$ and $Ax_j$ become linearly dependent for
$i\neq j$, then zero is an eigenvalue. Use this to find the remaining
eigenvalue of the matrix.}

\answer Linear dependence means that some scalar $c$ satisfies
$Ax_i=cAx_j$. By linearity,
\[
  A(x_i-cx_j)=Ax_i-cAx_j=0.
\]
The subspace-iteration state vectors are orthogonal, hence independent.
Therefore $x_i-cx_j\neq0$.

\conclusion{$A$ sends a nonzero vector to zero, so that vector is an
eigenvector with eigenvalue $0$.}

For the given matrix, dependence occurs in the first round:
\[
  Ae_2=(1,0,0),
  \qquad
  Ae_3=(2,0,0)=2Ae_2.
\]
Consequently,
\[
  A(2e_2-e_3)=0,
\]
and the remaining eigenvalue is $0$. The complete spectrum is
\[
  \boxed{\{\sqrt5,-\sqrt5,0\}}.
\]

\examlabel{Exam 3 --- Question 3}
\question{The two-bit addressing function is
\[
 \operatorname{addr}(x,y,z_{++},z_{+-},z_{-+},z_{--})=z_{xy}.
\]
\begin{enumerate}
\item Express it using the four point functions.
\item Calculate
$\widehat{\operatorname{addr}}(\varnothing)$,
$\widehat{\operatorname{addr}}(\{1,2,3\})$, and
$\widehat{\operatorname{addr}}(\{1,2,3,4\})$,
where variables $3$ and $4$ are $z_{++}$ and $z_{+-}$.
\end{enumerate}}

\answer For $a,b\in\{-1,+1\}$, define
\[
  \operatorname{point}_{ab}(x,y)
  =\frac{(1+ax)(1+by)}4.
\]
Exactly one point function equals $1$ for each input $(x,y)$, so
\[
\boxed{
  \operatorname{addr}
  =\sum_{a,b\in\{\pm1\}}
    z_{ab}\operatorname{point}_{ab}(x,y)}.
\]
Equivalently, this is the sum of the $++,+-,-+,--$ terms.

\textbf{Constant coefficient.}
Every summand contains one $z$ variable, so no constant term appears:
\[
  \boxed{\widehat{\operatorname{addr}}(\varnothing)=0}.
\]

\textbf{Coefficient indexed by $\{1,2,3\}$.}
This is the coefficient of $xyz_{++}$. Expand only the $++$ summand:
\[
\begin{aligned}
 z_{++}\operatorname{point}_{++}(x,y)
 &=\frac{z_{++}(1+x)(1+y)}4\\
 &=\frac14\left(
 z_{++}+xz_{++}+yz_{++}+xyz_{++}\right).
\end{aligned}
\]
Thus
\[
  \boxed{\widehat{\operatorname{addr}}(\{1,2,3\})=\frac14}.
\]

\textbf{Coefficient indexed by $\{1,2,3,4\}$.}
This asks for the monomial $xyz_{++}z_{+-}$. Every summand contains
exactly one $z$ variable, never a product of two different $z$ variables.
Therefore
\[
  \boxed{\widehat{\operatorname{addr}}(\{1,2,3,4\})=0}.
\]

\conclusion{Equivalently, every point function has degree $2$ in $x,y$,
and multiplying by one addressed bit gives degree at most $3$.
The degree-$4$ monomial $xyz_{++}z_{+-}$ cannot occur.}

\end{multicols*}

\end{document}
